\documentclass[10pt,journal,compsoc]{IEEEtran}

\usepackage{cite}
\usepackage{amsmath,amssymb,amsfonts}
\usepackage{algorithmic}
\usepackage{graphicx}
\usepackage{textcomp}
\usepackage{xcolor}
\usepackage{booktabs}
\usepackage{hyperref}

\begin{document}

\title{Multi-Scale Overlapping Pattern Tokenization (MSOPT) for Non-Stationary Financial Time Series}

\author{Suryaansh~Prithvijit~Singh,~\IEEEmembership{Lead Researcher,}
        and~Shivaji~Pawar,~\IEEEmembership{Faculty Supervisor}
\IEEEcompsocitemizethanks{
\IEEEcompsocthanksitem S. P. Singh and S. Pawar are with the Department of Future Tech, Universal AI University.\protect\\
E-mail: suryaansh.singh@universalai.edu.in
}}

\markboth{IEEE Transactions on Pattern Analysis and Machine Intelligence}%
{Singh \MakeLowercase{\textit{et al.}}: MSOPT for Non-Stationary Financial Time Series}

\IEEEtitleabstractindextext{%
\begin{abstract}
Time series modeling in non-stationary financial domains faces severe challenges due to low signal-to-noise ratios, regime shifts, and multi-scale visual pattern structures. Standard serial pointwise models and uniform non-overlapping patch Transformers (e.g., PatchTST) suffer from single-scale rigidity and token boundary clipping. In this paper, we propose Multi-Scale Overlapping Pattern Tokenization (MSOPT)—a framework that converts 1D price series into a 2D Scale-Time Spatial Tensor using dense dilated receptive field extraction ($s=1$) and 1D-SAX symbolic discretization. Empirical experiments under walk-forward evaluation protocols measure performance across equity and fixed-income assets post transaction costs.
\end{abstract}

\begin{IEEEkeywords}
Time Series Tokenization, Multi-Scale Pattern Recognition, 1D-SAX Discretization, Non-Stationary Financial Learning, Walk-Forward Backtesting.
\end{IEEEkeywords}}

\maketitle

\section{Introduction}
\IEEEPARstart{F}{inancial} time series analysis represents one of the most challenging domains in quantitative machine learning due to extreme non-stationarity, low signal-to-noise ratios (SNR), and non-linear regime shifts. Traditional sequence models process price data as a 1D sequence of continuous scalar points $x_t \in \mathbb{R}$. This serial approach exhibits quadratic computational complexity $\mathcal{O}(T^2)$ over long temporal horizons and fails to capture localized multi-scale visual shapes (e.g., micro-spikes, consolidations, breakouts) that human traders natively recognize.

Recent architectures, such as PatchTST \cite{nie2023patchtst}, aggregate scalar values into contiguous uniform patches ($P=16$). However, uniform non-overlapping patching forces rigid partition boundaries that clip pattern inflections arbitrarily. Furthermore, non-overlapping patch boundaries lack translation invariance.

To overcome these barriers, we propose \textbf{Multi-Scale Overlapping Pattern Tokenization (MSOPT)}. MSOPT treats financial series as 2D spatial chart primitives.

\section{Related Work}
Time series representation learning spans three dominant paradigms:
\begin{enumerate}
    \item \textbf{Pointwise Serial Architectures}: LSTMs, DeepAR, and N-BEATS process step-by-step inputs.
    \item \textbf{Fixed Uniform Patching}: PatchTST \cite{nie2023patchtst} aggregates subseries into non-overlapping patches.
    \item \textbf{Symbolic Discretization}: BORF \cite{spinnato2024borf} and SAX \cite{lin2003sax} discretize subseries into symbolic representations. TS-BPE \cite{gotz2025tsbpe} applies subseries Byte Pair Encoding.
\end{enumerate}

\section{Methodology}
\subsection{Multi-Scale Dilated Receptive Fields}
For receptive field window sizes $\mathcal{W} = \{4, 8, 16, 32\}$ and dilations $\mathcal{D} = \{1, 2, 4\}$, we sample subseries densely with stride $s=1$:
\begin{equation}
\mathbf{s}_{t}^{(w,d)} = \left[ x_{t - d \cdot (w-1)}, x_{t - d \cdot (w-2)}, \dots, x_{t} \right]
\end{equation}

\subsection{1D-SAX Symbolic Discretization}
Subseries are standardized and quantized into segment mean ($\alpha_{\mu}$) and slope ($\alpha_{\beta}$) symbols \cite{lin2003sax}, forming discrete 1D-SAX word strings.

\section{Empirical Evaluation}
We evaluate MSOPT under a 10-year expanding walk-forward protocol (2016--2025) across \texttt{SPY}, \texttt{QQQ}, \texttt{AAPL}, and \texttt{TLT}. All backtests enforce \textbf{5 bps (0.05\%) transaction costs} per trade flip.

\begin{table}[ht]
\caption{Master 10-Year Out-of-Sample Walk-Forward Performance (2016--2025, Post 5 Bps Costs)}
\label{tab:master_results}
\centering
\small
\begin{tabular}{llcccc}
\toprule
\textbf{Asset} & \textbf{Model Paradigm} & \textbf{OOS Acc} & \textbf{Net Return} & \textbf{Sharpe} & \textbf{Max DD} \\
\midrule
\texttt{SPY} & Baseline (Tech Lags) & 43.84\% & -46.46\% & -0.4043 & -56.72\% \\
\texttt{SPY} & \textbf{MSOPT Tokens (Ours)} & \textbf{45.68\%} & \textbf{+229.60\%} & \textbf{0.7589} & \textbf{-35.75\%} \\
\midrule
\texttt{QQQ} & Baseline (Tech Lags) & 44.60\% & +36.74\% & 0.2670 & -59.52\% \\
\texttt{QQQ} & \textbf{MSOPT Tokens (Ours)} & \textbf{44.88\%} & \textbf{+297.54\%} & \textbf{0.7341} & \textbf{-30.54\%} \\
\midrule
\texttt{AAPL} & Baseline (Tech Lags) & 43.20\% & +2354.43\% & 1.8004 & -23.57\% \\
\texttt{AAPL} & \textbf{MSOPT Tokens (Ours)} & \textbf{37.39\%} & \textbf{-27.66\%} & \textbf{0.0315} & \textbf{-56.33\%} \\
\midrule
\texttt{TLT} & Baseline (Tech Lags) & 40.22\% & -99.79\% & -4.7356 & -99.80\% \\
\texttt{TLT} & \textbf{MSOPT Tokens (Ours)} & \textbf{37.31\%} & \textbf{-30.64\%} & \textbf{-0.1757} & \textbf{-55.79\%} \\
\bottomrule
\end{tabular}
\end{table}

\subsection{Discussion of Findings}
On broad equity index ETFs (\texttt{SPY} and \texttt{QQQ}), MSOPT pattern tokens filter high-frequency noise and prevent excessive position churn. For \texttt{SPY}, MSOPT reduced total position flips over 10 years from 633 down to 71, improving Sharpe ratio from -0.4043 to +0.7589. On single-stock mega-cap momentum (\texttt{AAPL}), persistent trend-following baselines outperformed token-frequency features, highlighting that local shape words are most effective when applied to macro regime filtering.

\bibliographystyle{IEEEtran}
\bibliography{references}

\end{document}
